% ==================== THEME 5 — BOUCLES IMBRIQUEES (entraînement) ====================

% ---- imbr-trace-compteur ----
\begin{question}{ent-imbr-trace-compteur}
    Soit le code ci-dessous, qu'affichera le programme en fin d'exécution ?
    \lstinputlisting{sources/codes/ent-boucle-imbriquee.c}
    \begin{multicols}{3}
        \begin{reponses}
            \mauvaise{4}
            \mauvaise{3}
            \bonne{\checkmark\ 12}
            \mauvaise{7}
            \mauvaise{24}
            \mauvaise{Boucle infinie}
        \end{reponses}
    \end{multicols}
\end{question}

% ---- imbr-nb-tours ----
\begin{question}{ent-imbr-nb-tours}
    Combien de fois le corps de la boucle interne s'exécute-t-il au total dans \lstinline[language=C]|for (int i=0;i<4;i++) { for (int j=0;j<=i;j++) { } }| ?
    \begin{multicols}{5}
        \begin{reponses}
            \mauvaise{16}
            \mauvaise{4}
            \mauvaise{6}
            \bonne{\checkmark\ 10}
            \mauvaise{1}
        \end{reponses}
    \end{multicols}
\end{question}

% ---- imbr-variable-confusion ----
\begin{question}{ent-imbr-variable-confusion}
    Quel est l'effet de modifier la variable \texttt{i} à l'intérieur du corps de la boucle interne, si \texttt{i} est aussi la variable de la boucle externe ?
    \begin{multicols}{2}
        \begin{reponses}
            \bonne{\checkmark\ Cela perturbe le comptage de la boucle externe, qui utilise la même variable}
            \mauvaise{Aucun effet, chaque boucle possède sa propre copie de \texttt{i}}
            \mauvaise{Cela provoque une erreur de compilation}
            \mauvaise{Cela n'affecte que la boucle interne}
        \end{reponses}
    \end{multicols}
\end{question}

% ---- imbr-trace-variable (tableau structuré) ----
\begin{question}{ent-imbr-trace-variable}
    Pour chaque tour de la boucle externe, indiquer la valeur de \texttt{total} en fin de tour.
    \lstinputlisting[language=c]{sources/codes/ent-trace-imbriquee.c}
    \evaluationProfTableauTrace{2}{%
        \begin{center}
        \begin{tabular}{|c|p{4cm}|}
            \hline
            \textbf{Tour (valeur de \texttt{i})} & \textbf{Valeur de \texttt{total} en fin de tour} \\
            \hline
            1 & \textbf{8} \\[0.3cm]
            \hline
            2 & \textbf{6} \\[0.3cm]
            \hline
            3 & \textbf{4} \\[0.3cm]
            \hline
        \end{tabular}
        \end{center}
    }
\end{question}
